\documentclass[10pt,a4paper]{amsart}
\usepackage[margin=19mm]{geometry}
\usepackage{amsmath,amssymb,mathtools}
\usepackage[scr=rsfs,cal=euler]{mathalfa}
\usepackage{xfrac}
\usepackage[unicode,hidelinks,bookmarks=false]{hyperref}
\newcommand{\ie}{i.e.\ }
\newcommand{\cf}{cf.\ }
\newcommand{\resp}{respectively\ }
\newcommand{\Subsection}{\subsection}
\providecommand{\nobreakdash}{\nobreak\hbox{-}}
\setlength{\emergencystretch}{2em}
\multlinegap0pt
\usepackage{amssymb}
\usepackage{float}
\usepackage{mathtools}
\usepackage{enumerate}
\usepackage{xcolor}
\usepackage{tikz}
\newtheorem{proposition}{Proposition}
\title{An Alternating Primal Heuristic for Nonconvex MIQCQP with Dynamic Convexification and Parallel Local Branching}
\author{Yongzheng Dai, Chen Chen}
\date{Source: arXiv:2604.04417v1 (CC BY 4.0)}
\begin{document}
\maketitle

\noindent\emph{Source and licence.} An Alternating Primal Heuristic for Nonconvex MIQCQP with Dynamic Convexification and Parallel Local Branching. Version arXiv:2604.04417v1 (CC BY 4.0), distributed by arXiv under the Creative Commons Attribution 4.0 International licence, http://creativecommons.org/licenses/by/4.0/, as recorded in the arXiv licence metadata for this exact version and shown on the abstract page https://arxiv.org/abs/2604.04417v1. Full licence notice: \texttt{licenses/arxiv-2604.04417-CC-BY-4.0.txt}.\par\smallskip

\noindent\textit{Editorial scope.} Counted unit: the article's proposition prop:fixed-point (source0.tex lines 258--268). The article's complete original proof of it is delivered with it (source0.tex lines 269--301, one \texttt{proof} environment of 33 lines). No mathematics is written, repaired or re-derived: the statement and the proof are the article's own lines, in the article's own notation, and no retained line is substituted or re-flowed.  No further source-local context is needed: every cross-reference of every retained line resolves inside the two delivered spans.  Nothing was omitted from any retained span, so the package carries no omission record.  No retained line contains a citation, so the wrapper carries no bibliography and no bibliography entry.  No retained line contains a figure environment, an image-inclusion command or drawing code, so no image is delivered and the package declares no asset dependency.  Editorial formatting only, in addition: an AMS \texttt{amsart} preamble that restates the article's own declarations and macros used by the retained lines, namely the environments \texttt{proposition} on separate counters, together with the standard packages those lines need.  The retained environments print in this document as Proposition 1 in order of appearance, whereas the article prints the same items with the numbering of its own full document.  The complete native source of the article as distributed in the arXiv e-print for this exact version is bundled unchanged as \texttt{sources/197880/source0.tex}.
\begin{proposition}\label{prop:fixed-point}
    Define $\mathrm{Q}_{\mathrm{MIQP}}:\min_x\ x^TQ^0x + a^Tx$ with $n_\alpha$ binary variables $x_\alpha$, $n_\beta$ continuous variables $x_\beta \geq 0$, and
    \[Q^0=\begin{bmatrix}
            Q_\alpha\ Q_\gamma \\
            Q_\gamma^T\ Q_\beta
        \end{bmatrix}, a = \begin{bmatrix}
            a_\alpha\\
            a_\beta
        \end{bmatrix}.\]
    Let $x(\hat{u})=(x_\alpha(\hat{u}), x_\beta(\hat{u}))$ be an optimal solution to $\mathrm{Approx}(\hat{u})$ for $\mathrm{Q}_{\mathrm{MIQP}}$. If $\hat{u}_\beta=2x_\beta(\hat{u})$, then $x_\beta(\hat{u})$ satisfies stationarity of $\mathrm{Q}_{\mathrm{MIQP}}$ for fixed binary variables $x_\alpha(\hat{u})$.
\end{proposition}

\begin{proof}
    Let
    \[f(x) := x^TQ^0x + a^Tx = x_\alpha^T Q_\alpha x_\alpha + 2x_\alpha^T Q_\gamma x_\beta + x_\beta^T Q_\beta x_\beta + a_\alpha^Tx_\alpha+a_\beta^Tx_\beta.\]
    Suppose $x^*=(x^*_\alpha,x^*_\beta)$ is an optimal solution to $\mathrm{Q}_{\mathrm{MIQP}}$. Fixing $x_\alpha = x_\alpha^*$, the first-order stationarity condition of $x_\beta^* = \arg\min\{f(x):x_\alpha = x_\alpha^*, x_\beta\geq 0\}$ is
    \begin{equation}\label{eq:opt_condition}
        \left\{\begin{aligned}
        \nabla_{(x_\beta)_i} f(x^*) = 0, &\mbox{ if } (x_\beta^*)_i> 0\\
        \nabla_{(x_\beta)_i} f(x^*) \geq 0, &\mbox{ if } (x_\beta^*)_i= 0
    \end{aligned} \right.
    \end{equation}
    where
    \begin{equation*}
       \nabla_{x_\beta} f(x^*) =  2Q_\beta x^*_\beta + 2Q_\gamma^T x_\alpha^* +a_\beta.
    \end{equation*}

    Now, let us write $\mbox{Approx}(\hat{u}):=\arg\min\{g(x, \hat{u})\}$, where \[g(x, \hat{u}):=x^T(Q^0-\lambda_{\min}^0 I)x+a^Tx+\lambda^0_{\min}\hat{u}_\beta^Tx_\beta + \lambda_{\min}^0\mathbf{1}^T x_\alpha.\] Fixing $x_{\alpha} = x_\alpha(\hat{u})$, $g(x,\hat{u})$ is convex from $Q^0-\lambda_{\min}^0 I \succeq 0$, and we have the first-order optimality condition
    \[\left\{\begin{aligned}
        \nabla_{(x_\beta)_i} g(x(\hat{u}), \hat{u}) = 0, &\mbox{ if } (x_\beta(\hat{u}))_i> 0\\
        \nabla_{(x_\beta)_i} g(x(\hat{u}), \hat{u}) \geq 0, &\mbox{ if } (x_\beta(\hat{u}))_i= 0
    \end{aligned} \right.,\]
    where
    \[\nabla_{x_\beta}g(x(\hat{u}), \hat{u}) = 2(Q_\beta-\lambda_{\min}^0 I_\beta)x_\beta(\hat{u}) + 2Q_\gamma^T x_\alpha(\hat{u}) + a_\beta +\lambda_{\min}^0\hat{u}_\beta.\]
    Plugging in the relation $\hat{u}_\beta=2x_\beta(\hat{u})$,
    \begin{equation*}
        \begin{aligned}
            \nabla_{x_\beta}g(x(\hat{u}), \hat{u}) &= 2(Q_\beta-\lambda_{\min}^0 I_\beta)x_\beta(\hat{u}) +2Q_\gamma^T x_\alpha(\hat{u}) + a_\beta +2\lambda_{\min}^0 x_\beta(\hat{u})\\
            &=2Q_\beta x_\beta(\hat{u})+2Q_\gamma^T x_\alpha(\hat{u}) + a_\beta\\
            &=\nabla_{x_\beta} f(x(\hat{u})).
        \end{aligned}
    \end{equation*}
    which implies $(x_\alpha(\hat{u}), x_\beta(\hat{u}))$ satisfies the stationarity condition~(\ref{eq:opt_condition}).
    \qed
\end{proof}

\end{document}
